problem:  
Let \( (M^n,g) \) be a compact, simply connected Riemannian manifold with strictly positive sectional curvature.  
At each \(p\in M\), define the *directional pinching ratio*  
\[
\delta_v(p)=\frac{\min\{K_p(\sigma): v\in \sigma\subset T_pM\}}{\max\{K_p(\sigma): v\in \sigma\subset T_pM\}},\qquad v\in S_pM.
\]
Define the *integrated directional pinching constant*  
\[
\overline{\delta}(M)=\frac{1}{\operatorname{Vol}(SM)}\int_{SM}\delta_v(p)\,d\mu(p,v),
\]
where \(SM\) is the unit tangent bundle with the Sasaki measure.  

**Problem:**  
Assume that \(\overline{\delta}(M)>\tfrac{1}{4}\) and that sectional curvature attains values both above and below its median (so \(K\) is not constant).   
Is it necessarily true that \(M\) is diffeomorphic to a compact rank‑one symmetric space (CROSS)?  
Equivalently, can there exist a sequence of simply connected, positively curved manifolds \((M_i^n,g_i)\) with  
\[
\overline{\delta}(M_i)\to\tfrac{1}{4}^+,
\]
yet whose smooth types are not among those of CROSSes?  
If such a sequence exists, does any subsequence \( (M_{i_j},g_{i_j}) \) exhibit geometric degeneration (collapse, formation of singular limits, or concentration of curvature along a subset of directions)?

Why is it a "good" problem:  
This problem introduces a *directionally averaged* curvature pinching invariant, bridging **microscopic directional geometry** and **global topology**. Unlike classical pointwise or global pinching constants, the functional \(\overline{\delta}(M)\) measures how curvature variability distributes across directions rather than points—capturing a deep geometric–statistical property of curved manifolds.  

The problem asks whether *average directional rigidity* above the quarter‑pinching threshold enforces topological rigidity, or whether directional variance allows new flexible positively curved structures to emerge. This formulation intertwines **Riemannian geometry**, **integral geometry** (through averaging on the unit tangent bundle), and **topological classification** of smooth manifolds with positive curvature.  

It is conceptually minimal—expressed with an elegant integral ratio—but potentially profound: progress would require developing new tools relating averaged curvature distributions to topological obstructions, possibly advancing Ricci‑flow stability or Cheeger‑Gromov convergence theory. The question simultaneously generalizes and refines the classical pinching paradigm, making it a genuinely *good* and forward‑looking research problem.